% 双线性 phi(x,y)=xy 的一步更新，F(u)=(y,-x)，步长 eta=1/2，旧点 u=(1.6,0)。
% 普通一步：u - eta F(u) = (1.6,0.8)（亦为外梯度的试探点）；再走一步：(1.2,1.6)，半径 2；
% 外梯度校正：u - eta F(ubar) = (1.2,0.8)，半径 sqrt(2.08)≈1.44 < 1.6。
\begin{tikzpicture}[scale=1.55,note/.style={font=\FigureNoteFont\fontsize{7.5}{9}\selectfont,align=left},
  pt/.style={circle,fill=#1,inner sep=0pt,minimum size=3.6pt}]
  % rotation field -F(u)=(-y,x)
  \foreach \a in {-0.4,0.4,1.2,2.0}{\foreach \b in {-0.4,0.4,1.2,2.0}{
    \draw[->,FigureGrid,line width=.6pt] (\a,\b) -- ++({-0.16*\b},{0.16*\a});}}
  \draw[->,FigureStroke,line width=.6pt] (-0.6,0)--(2.4,0) node[right]{$x$};
  \draw[->,FigureStroke,line width=.6pt] (0,-0.6)--(0,2.3) node[above]{$y$};
  \draw[FigureMuted,dashed,line width=.6pt] (1.6,0) arc[start angle=0,end angle=80,radius=1.6];
  \node[pt=FigureInk,minimum size=4pt] at (0,0) {};
  \node[below left=0pt] at (0,0) {鞍点};
  \coordinate (u) at (1.6,0); \coordinate (ub) at (1.6,0.8);
  \coordinate (g2) at (1.2,1.6); \coordinate (eg) at (1.2,0.8);
  \draw[-{Stealth[length=5pt]},FigureBlue,line width=.9pt,shorten >=2pt] (u)--(ub);
  \draw[-{Stealth[length=5pt]},FigureBlue,line width=.9pt,dashed,shorten >=2pt] (ub)--(g2);
  \draw[-{Stealth[length=5pt]},FigureGreen,line width=1.1pt,shorten >=2pt] (u)--(eg);
  \node[pt=FigureInk] at (u) {}; \node[pt=FigureBlue] at (ub) {};
  \node[pt=FigureBlue] at (g2) {}; \node[pt=FigureGreen] at (eg) {};
  \node[below=1pt] at (u) {$\bm u$};
  \node[right=1pt,text=FigureBlue] at (ub) {$\bar{\bm u}$（试探点）};
  \node[right=1pt,text=FigureBlue] at (g2) {连续两次普通步};
  \node[left=1pt,text=FigureGreen] at (eg) {$\bm u^{+}$};
  \node[note,text=FigureMuted] at (0.95,-0.42) {虚线圆弧：与鞍点等距};
\end{tikzpicture}
